Rios são a principal via pela qual o plástico descartado em terra chega ao oceano. Se isso é verdade, a quantidade de plástico na foz de um rio deveria crescer com o número de pessoas que vivem na bacia que o alimenta. Esta dissertação testa essa hipótese na América do Norte e Central, dos Estados Unidos e sul do Canadá até o Canal do Panamá, cruzando três bases independentes: o relevo global GEBCO 2024, a grade de população GPW v4 (2020) e as amostras de \textit{pellets} plásticos do programa \textit{Nurdle Patrol} no banco de microplásticos marinhos da NOAA. As bacias hidrográficas são delineadas pelo algoritmo baseado em percolação invasiva (IPBA), tomando como saída cada pixel de terra vizinho ao oceano; as 138 maiores bacias, que cobrem 90\% da área de terra, formam a amostra. As cidades são definidas sem fronteiras administrativas pelo \textit{City Clustering Algorithm} (CCA), cujos dois parâmetros, o limiar de densidade $D^{*}$ e a distância de ligação $\ell$, foram calibrados pela lei de Zipf numa varredura de $300 \times 50$ combinações, resolvida em tempo quase linear por uma estrutura de união e busca que ativa os pixels em ordem decrescente de densidade. O critério escolhido, que exige expoente próximo de 1 simultaneamente para as 50, 100 e 200 maiores manchas, resultou em $D^{*} = 1010$ hab/km$^2$ e $\ell = 1{,}7$ km, com 12.185 manchas urbanas e expoentes de Hill entre 0,93 e 1,08. Manchas que cruzam divisores de água são repartidas entre as bacias. Cada bacia recebe o ponto de coleta mais próximo da sua foz, e o raio de aceitação é varrido de 1 em 1 km até 50 km. Com o raio de 38 km, que reúne a maioria dos pontos de coleta, 52 bacias têm medição, 29 delas com contagem positiva. Nessas 29, a concentração de \textit{pellets} cresce de forma sublinear com a população, com expoente $\beta = 0{,}28$ (IC 95\%: 0,11 a 0,54; $r = 0{,}37$; $p = 0{,}05$), e o indicador de quantidade, concentração vezes área, cresce com $\beta = 0{,}85$ ($r = 0{,}66$). A associação é fraca e frágil: incluindo as bacias com contagem zero, a correlação de postos cai para $\rho_s = 0{,}25$ ($p = 0{,}07$), e boa parte do expoente do indicador de quantidade vem da correlação entre população e área das bacias. Com raio de 1 km apenas três bacias têm medição. Conclui-se que os dados públicos atuais sustentam, no máximo, uma relação sublinear e fraca entre população da bacia e microplástico na foz, e que a principal limitação é a cobertura espacial das coletas, concentrada no Golfo do México. A dissertação contribui com um \textit{pipeline} reprodutível que integra percolação invasiva, agrupamento urbano calibrado por Zipf e alometria por bacia.

\palavraschave{microplástico; bacias hidrográficas; percolação invasiva; lei de Zipf; alometria urbana.}
